%!TeX root=../persuasiontop.tex
\chapter[Chapter \thechapter]{}

\begin{figure}[t!]
\centering
\includegraphics[width=\linewidth]{nouveau17}
\end{figure}

\lettrine[lraise=0.3]{W}{hile} Sir Walter and Elizabeth were assiduously pushing their good fortune in Laura Place, Anne was renewing an acquaintance of a very different description.

She had called on her former governess, and had heard from her of there being an old schoolfellow in Bath, who had the two strong claims on her attention of past kindness and present suffering. Miss Hamilton, now Mrs Smith, had shewn her kindness in one of those periods of her life when it had been most valuable. Anne had gone unhappy to school, grieving for the loss of a mother whom she had dearly loved, feeling her separation from home, and suffering as a girl of fourteen, of strong sensibility and not high spirits, must suffer at such a time; and Miss Hamilton, three years older than herself, but still from the want of near relations and a settled home, remaining another year at school, had been useful and good to her in a way which had considerably lessened her misery, and could never be remembered with indifference.

Miss Hamilton had left school, had married not long afterwards, was said to have married a man of fortune, and this was all that Anne had known of her, till now that their governess's account brought her situation forward in a more decided but very different form.

She was a widow and poor. Her husband had been extravagant; and at his death, about two years before, had left his affairs dreadfully involved. She had had difficulties of every sort to contend with, and in addition to these distresses had been afflicted with a severe rheumatic fever, which, finally settling in her legs, had made her for the present a cripple. She had come to Bath on that account, and was now in lodgings near the hot baths, living in a very humble way, unable even to afford herself the comfort of a servant, and of course almost excluded from society.

Their mutual friend answered for the satisfaction which a visit from Miss Elliot would give Mrs Smith, and Anne therefore lost no time in going. She mentioned nothing of what she had heard, or what she intended, at home. It would excite no proper interest there. She only consulted Lady Russell, who entered thoroughly into her sentiments, and was most happy to convey her as near to Mrs Smith's lodgings in Westgate Buildings, as Anne chose to be taken.

The visit was paid, their acquaintance re-established, their interest in each other more than re-kindled. The first ten minutes had its awkwardness and its emotion. Twelve years were gone since they had parted, and each presented a somewhat different person from what the other had imagined. Twelve years had changed Anne from the blooming, silent, unformed girl of fifteen, to the elegant little woman of seven-and-twenty, with every beauty except bloom, and with manners as consciously right as they were invariably gentle; and twelve years had transformed the fine-looking, well-grown Miss Hamilton, in all the glow of health and confidence of superiority, into a poor, infirm, helpless widow, receiving the visit of her former protegee as a favour; but all that was uncomfortable in the meeting had soon passed away, and left only the interesting charm of remembering former partialities and talking over old times.

Anne found in Mrs Smith the good sense and agreeable manners which she had almost ventured to depend on, and a disposition to converse and be cheerful beyond her expectation. Neither the dissipations of the past—and she had lived very much in the world—nor the restrictions of the present, neither sickness nor sorrow seemed to have closed her heart or ruined her spirits.

In the course of a second visit she talked with great openness, and Anne's astonishment increased. She could scarcely imagine a more cheerless situation in itself than Mrs Smith's. She had been very fond of her husband: she had buried him. She had been used to affluence: it was gone. She had no child to connect her with life and happiness again, no relations to assist in the arrangement of perplexed affairs, no health to make all the rest supportable. Her accommodations were limited to a noisy parlour, and a dark bedroom behind, with no possibility of moving from one to the other without assistance, which there was only one servant in the house to afford, and she never quitted the house but to be conveyed into the warm bath. Yet, in spite of all this, Anne had reason to believe that she had moments only of languor and depression, to hours of occupation and enjoyment. How could it be? She watched, observed, reflected, and finally determined that this was not a case of fortitude or of resignation only. A submissive spirit might be patient, a strong understanding would supply resolution, but here was something more; here was that elasticity of mind, that disposition to be comforted, that power of turning readily from evil to good, and of finding employment which carried her out of herself, which was from nature alone. It was the choicest gift of Heaven; and Anne viewed her friend as one of those instances in which, by a merciful appointment, it seems designed to counterbalance almost every other want.

There had been a time, Mrs Smith told her, when her spirits had nearly failed. She could not call herself an invalid now, compared with her state on first reaching Bath. Then she had, indeed, been a pitiable object; for she had caught cold on the journey, and had hardly taken possession of her lodgings before she was again confined to her bed and suffering under severe and constant pain; and all this among strangers, with the absolute necessity of having a regular nurse, and finances at that moment particularly unfit to meet any extraordinary expense. She had weathered it, however, and could truly say that it had done her good. It had increased her comforts by making her feel herself to be in good hands. She had seen too much of the world, to expect sudden or disinterested attachment anywhere, but her illness had proved to her that her landlady had a character to preserve, and would not use her ill; and she had been particularly fortunate in her nurse, as a sister of her landlady, a nurse by profession, and who had always a home in that house when unemployed, chanced to be at liberty just in time to attend her. <And she,> said Mrs Smith, <besides nursing me most admirably, has really proved an invaluable acquaintance. As soon as I could use my hands she taught me to knit, which has been a great amusement; and she put me in the way of making these little thread-cases, pin-cushions and card-racks, which you always find me so busy about, and which supply me with the means of doing a little good to one or two very poor families in this neighbourhood. She had a large acquaintance, of course professionally, among those who can afford to buy, and she disposes of my merchandise. She always takes the right time for applying. Everybody's heart is open, you know, when they have recently escaped from severe pain, or are recovering the blessing of health, and Nurse Rooke thoroughly understands when to speak. She is a shrewd, intelligent, sensible woman. Hers is a line for seeing human nature; and she has a fund of good sense and observation, which, as a companion, make her infinitely superior to thousands of those who having only received <the best education in the world,> know nothing worth attending to. Call it gossip, if you will, but when Nurse Rooke has half an hour's leisure to bestow on me, she is sure to have something to relate that is entertaining and profitable: something that makes one know one's species better. One likes to hear what is going on, to be \textit{au fait} as to the newest modes of being trifling and silly. To me, who live so much alone, her conversation, I assure you, is a treat.>

Anne, far from wishing to cavil at the pleasure, replied, <I can easily believe it. Women of that class have great opportunities, and if they are intelligent may be well worth listening to. Such varieties of human nature as they are in the habit of witnessing! And it is not merely in its follies, that they are well read; for they see it occasionally under every circumstance that can be most interesting or affecting. What instances must pass before them of ardent, disinterested, self-denying attachment, of heroism, fortitude, patience, resignation: of all the conflicts and all the sacrifices that ennoble us most. A sick chamber may often furnish the worth of volumes.>

<Yes,> said Mrs Smith more doubtingly, <sometimes it may, though I fear its lessons are not often in the elevated style you describe. Here and there, human nature may be great in times of trial; but generally speaking, it is its weakness and not its strength that appears in a sick chamber: it is selfishness and impatience rather than generosity and fortitude, that one hears of. There is so little real friendship in the world! and unfortunately> (speaking low and tremulously) <there are so many who forget to think seriously till it is almost too late.>

Anne saw the misery of such feelings. The husband had not been what he ought, and the wife had been led among that part of mankind which made her think worse of the world than she hoped it deserved. It was but a passing emotion however with Mrs Smith; she shook it off, and soon added in a different tone—

<I do not suppose the situation my friend Mrs Rooke is in at present, will furnish much either to interest or edify me. She is only nursing Mrs Wallis of Marlborough Buildings; a mere pretty, silly, expensive, fashionable woman, I believe; and of course will have nothing to report but of lace and finery. I mean to make my profit of Mrs Wallis, however. She has plenty of money, and I intend she shall buy all the high-priced things I have in hand now.>

Anne had called several times on her friend, before the existence of such a person was known in Camden Place. At last, it became necessary to speak of her. Sir Walter, Elizabeth and Mrs Clay, returned one morning from Laura Place, with a sudden invitation from Lady Dalrymple for the same evening, and Anne was already engaged, to spend that evening in Westgate Buildings. She was not sorry for the excuse. They were only asked, she was sure, because Lady Dalrymple being kept at home by a bad cold, was glad to make use of the relationship which had been so pressed on her; and she declined on her own account with great alacrity—<She was engaged to spend the evening with an old schoolfellow.> They were not much interested in anything relative to Anne; but still there were questions enough asked, to make it understood what this old schoolfellow was; and Elizabeth was disdainful, and Sir Walter severe.

<Westgate Buildings!> said he, <and who is Miss Anne Elliot to be visiting in Westgate Buildings? A Mrs Smith. A widow Mrs Smith; and who was her husband? One of five thousand Mr Smiths whose names are to be met with everywhere. And what is her attraction? That she is old and sickly. Upon my word, Miss Anne Elliot, you have the most extraordinary taste! Everything that revolts other people, low company, paltry rooms, foul air, disgusting associations are inviting to you. But surely you may put off this old lady till to-morrow: she is not so near her end, I presume, but that she may hope to see another day. What is her age? Forty?>

<No, sir, she is not one-and-thirty; but I do not think I can put off my engagement, because it is the only evening for some time which will at once suit her and myself. She goes into the warm bath to-morrow, and for the rest of the week, you know, we are engaged.>

<But what does Lady Russell think of this acquaintance?> asked Elizabeth.

<She sees nothing to blame in it,> replied Anne; <on the contrary, she approves it, and has generally taken me when I have called on Mrs Smith.>

<Westgate Buildings must have been rather surprised by the appearance of a carriage drawn up near its pavement,> observed Sir Walter. <Sir Henry Russell's widow, indeed, has no honours to distinguish her arms, but still it is a handsome equipage, and no doubt is well known to convey a Miss Elliot. A widow Mrs Smith lodging in Westgate Buildings! A poor widow barely able to live, between thirty and forty; a mere Mrs Smith, an every-day Mrs Smith, of all people and all names in the world, to be the chosen friend of Miss Anne Elliot, and to be preferred by her to her own family connections among the nobility of England and Ireland! Mrs Smith! Such a name!>

Mrs Clay, who had been present while all this passed, now thought it advisable to leave the room, and Anne could have said much, and did long to say a little in defence of \textit{her} friend's not very dissimilar claims to theirs, but her sense of personal respect to her father prevented her. She made no reply. She left it to himself to recollect, that Mrs Smith was not the only widow in Bath between thirty and forty, with little to live on, and no surname of dignity.

Anne kept her appointment; the others kept theirs, and of course she heard the next morning that they had had a delightful evening. She had been the only one of the set absent, for Sir Walter and Elizabeth had not only been quite at her ladyship's service themselves, but had actually been happy to be employed by her in collecting others, and had been at the trouble of inviting both Lady Russell and Mr Elliot; and Mr Elliot had made a point of leaving Colonel Wallis early, and Lady Russell had fresh arranged all her evening engagements in order to wait on her. Anne had the whole history of all that such an evening could supply from Lady Russell. To her, its greatest interest must be, in having been very much talked of between her friend and Mr Elliot; in having been wished for, regretted, and at the same time honoured for staying away in such a cause. Her kind, compassionate visits to this old schoolfellow, sick and reduced, seemed to have quite delighted Mr Elliot. He thought her a most extraordinary young woman; in her temper, manners, mind, a model of female excellence. He could meet even Lady Russell in a discussion of her merits; and Anne could not be given to understand so much by her friend, could not know herself to be so highly rated by a sensible man, without many of those agreeable sensations which her friend meant to create.

Lady Russell was now perfectly decided in her opinion of Mr Elliot. She was as much convinced of his meaning to gain Anne in time as of his deserving her, and was beginning to calculate the number of weeks which would free him from all the remaining restraints of widowhood, and leave him at liberty to exert his most open powers of pleasing. She would not speak to Anne with half the certainty she felt on the subject, she would venture on little more than hints of what might be hereafter, of a possible attachment on his side, of the desirableness of the alliance, supposing such attachment to be real and returned. Anne heard her, and made no violent exclamations; she only smiled, blushed, and gently shook her head.

<I am no match-maker, as you well know,> said Lady Russell, <being much too well aware of the uncertainty of all human events and calculations. I only mean that if Mr Elliot should some time hence pay his addresses to you, and if you should be disposed to accept him, I think there would be every possibility of your being happy together. A most suitable connection everybody must consider it, but I think it might be a very happy one.>

<Mr Elliot is an exceedingly agreeable man, and in many respects I think highly of him,> said Anne; <but we should not suit.>

Lady Russell let this pass, and only said in rejoinder, <I own that to be able to regard you as the future mistress of Kellynch, the future Lady Elliot, to look forward and see you occupying your dear mother's place, succeeding to all her rights, and all her popularity, as well as to all her virtues, would be the highest possible gratification to me. You are your mother's self in countenance and disposition; and if I might be allowed to fancy you such as she was, in situation and name, and home, presiding and blessing in the same spot, and only superior to her in being more highly valued! My dearest Anne, it would give me more delight than is often felt at my time of life!>

Anne was obliged to turn away, to rise, to walk to a distant table, and, leaning there in pretended employment, try to subdue the feelings this picture excited. For a few moments her imagination and her heart were bewitched. The idea of becoming what her mother had been; of having the precious name of <Lady Elliot> first revived in herself; of being restored to Kellynch, calling it her home again, her home for ever, was a charm which she could not immediately resist. Lady Russell said not another word, willing to leave the matter to its own operation; and believing that, could Mr Elliot at that moment with propriety have spoken for himself!—she believed, in short, what Anne did not believe. The same image of Mr Elliot speaking for himself brought Anne to composure again. The charm of Kellynch and of <Lady Elliot> all faded away. She never could accept him. And it was not only that her feelings were still adverse to any man save one; her judgement, on a serious consideration of the possibilities of such a case, was against Mr Elliot.

Though they had now been acquainted a month, she could not be satisfied that she really knew his character. That he was a sensible man, an agreeable man, that he talked well, professed good opinions, seemed to judge properly and as a man of principle, this was all clear enough. He certainly knew what was right, nor could she fix on any one article of moral duty evidently transgressed; but yet she would have been afraid to answer for his conduct. She distrusted the past, if not the present. The names which occasionally dropt of former associates, the allusions to former practices and pursuits, suggested suspicions not favourable of what he had been. She saw that there had been bad habits; that Sunday travelling had been a common thing; that there had been a period of his life (and probably not a short one) when he had been, at least, careless in all serious matters; and, though he might now think very differently, who could answer for the true sentiments of a clever, cautious man, grown old enough to appreciate a fair character? How could it ever be ascertained that his mind was truly cleansed?

Mr Elliot was rational, discreet, polished, but he was not open. There was never any burst of feeling, any warmth of indignation or delight, at the evil or good of others. This, to Anne, was a decided imperfection. Her early impressions were incurable. She prized the frank, the open-hearted, the eager character beyond all others. Warmth and enthusiasm did captivate her still. She felt that she could so much more depend upon the sincerity of those who sometimes looked or said a careless or a hasty thing, than of those whose presence of mind never varied, whose tongue never slipped.

Mr Elliot was too generally agreeable. Various as were the tempers in her father's house, he pleased them all. He endured too well, stood too well with every body. He had spoken to her with some degree of openness of Mrs Clay; had appeared completely to see what Mrs Clay was about, and to hold her in contempt; and yet Mrs Clay found him as agreeable as any body.

Lady Russell saw either less or more than her young friend, for she saw nothing to excite distrust. She could not imagine a man more exactly what he ought to be than Mr Elliot; nor did she ever enjoy a sweeter feeling than the hope of seeing him receive the hand of her beloved Anne in Kellynch church, in the course of the following autumn.